% !TEX root = ../../../main.tex
% 所属：第一部分：单变量分析学 / 第7章：微分法
% 节：7.8 Taylor公式
% 正文按原文件顺序保存；由 ../chapter7.tex 调用。
\section{Taylor公式}
\begin{theorem}(带 Peano 余项的 Taylor 公式)
	设 $f(x)$ 在 $x = x_0$ 处有 $n$ 阶导数.则存在 $x_0$ 的邻域 $O(x_0, \delta)$，在 $O(x_0, \delta)$ 上，
	\[
	f(x) = f(x_0) + f'(x_0)(x-x_0) + \frac{1}{2!} f''(x_0)(x-x_0)^2 + \cdots + \frac{1}{n!} f^{(n)}(x_0)(x-x_0)^n + r_n(x)
	\]
	
	其中，$r_n(x) = o((x-x_0)^n)$，称为 Peano 余项.
	
	
\end{theorem}
\begin{proof}
	设
	\[
	r_n(x) = f(x) - \left[f(x_0) + f'(x_0)(x-x_0) + \cdots + \frac{1}{n!} f^{(n)}(x_0)(x-x_0)^n\right]
	\]
	下证 $r_n(x) = o((x-x_0)^n)$，也即
	\[
	\lim_{x \to x_0} \frac{r_n(x)}{(x-x_0)^n} = 0
	\]
	对 $n$ 用归纳法. $n=1$ 时，由导数的定义，
	\[
	f(x) = f(x_0) + f'(x_0)(x-x_0) + o(x-x_0),
	\]
	即 $r_1(x) = o(x-x_0)$，成立.
	
	假设结论对 $n$ 成立. 设 $f$ 在 $x_0$ 有 $n+1$ 阶导数，则 $f'$ 在 $x_0$ 有 $n$ 阶导数.
	将 $f$ 写成 $f(x)=P_{n+1}(x)+r_{n+1}(x)$，其中 $P_{n+1}$ 是 $n+1$ 次 Taylor 多项式.
	
	注意到 $r_{n+1}(x_0) = r_{n+1}'(x_0) = \cdots = r_{n+1}^{(n)}(x_0) = 0$，
	且 $r_{n+1}^{(n+1)}(x_0) = 0$. 反复应用 Cauchy 中值定理：
	\[
	\frac{r_{n+1}(x)}{(x-x_0)^{n+1}}
	= \frac{r_{n+1}'(\xi_1)}{(n+1)(\xi_1-x_0)^n}
	= \cdots
	= \frac{r_{n+1}^{(n)}(\xi_n)}{(n+1)!\,(\xi_n-x_0)},
	\]
	其中 $\xi_k$ 介于 $x_0$ 与 $\xi_{k-1}$ 之间（$\xi_0 = x$）.
	最后一步，
	\[
	\frac{r_{n+1}^{(n)}(\xi_n)}{\xi_n-x_0}
	= \frac{r_{n+1}^{(n)}(\xi_n) - r_{n+1}^{(n)}(x_0)}{\xi_n-x_0}
	\to r_{n+1}^{(n+1)}(x_0) = 0 \quad (\xi_n\to x_0).
	\]
	故 $r_{n+1}(x) = o((x-x_0)^{n+1})$，归纳完成.
\end{proof}
\begin{theorem}(带 Lagrange 余项的 Taylor 公式)
	设 $f(x) \in D^n[a, b]\cap D^{n+1}(a, b)$，任取 $ x_0 \in (a,b)$，$\exists \xi \in (a, b)$ 介于 $x$ 与 $x_0$ 之间，使得
	\[
	f(x) = f(x_0) + f'(x_0)(x-x_0) + \cdots + \frac{1}{n!} f^{(n)}(x_0)(x-x_0)^n + \frac{f^{(n+1)}(\xi)}{(n+1)!}(x-x_0)^{n+1}
	\]
	$r_n(x)=\dfrac{f^{(n+1)}(\xi)}{(n+1)!}(x-x_0)^{n+1}$称为Lagrange余项
\end{theorem}
\begin{proof}
	考虑辅助函数
	\[
	G(t) = f(x) - \sum_{k=0}^n \frac{1}{k!} f^{(k)}(t)(x-t)^k \quad \text{和} \quad H(t) = (x-t)^{n+1}.
	\]
	
	那么定理的结论（即需要证明的）是
	\[
	G(x_0) = \frac{f^{(n+1)}(\xi)}{(n+1)!} H(x_0).
	\]
	
	不妨设 $x_0 < x$（$x < x_0$ 时证明类似）.则 $G(t)$ 和 $H(t)$ 在 $[x_0, x]$ 上连续，在 $(x_0, x)$ 上可导，且
	
	\[\begin{aligned}
		G'(t) =& -[f'(t) + (-f'(t)) + f''(t)(x-t) + (-f''(t))(x-t)+ \frac{1}{2!} f'''(t)(x-t)^2 + \cdots\\&-\frac{1}{(n-1)!} f^{(n)}(t)(x-t)^{n-1} + \frac{f^{(n+1)}(t)}{n!}(x-t)^n]\\
		=&-\frac{1}{n!} f^{(n+1)}(t)(x-t)^n
	\end{aligned}\]
	
	\[
	H'(t) = -(n+1)(x-t)^n.
	\]
	
	显然 $H'(t)$ 在 $(x_0, x)$ 上不等于零.因为 $G(x) = H(x) = 0$，由 Cauchy 中值定理可得
	\[
	\frac{G(x_0)}{H(x_0)} = \frac{G(x) - G(x_0)}{H(x) - H(x_0)} = \frac{G'(\xi)}{H'(\xi)} = \frac{f^{(n+1)}(\xi)}{(n+1)!}, \quad \xi \in (x_0, x),
	\]
	
	由此 $G(x_0) = \dfrac{f^{(n+1)}(\xi)}{(n+1)!} H(x_0)$.
\end{proof}
\begin{corollary}(Maclaurin 公式)
	在 $x_0=0$ 时的 Taylor 公式，即为 Maclaurin 公式，
	$$
	f(x) = f(0) + f'(0)x + \frac{1}{2!} f''(0)x^2 + \cdots + \frac{1}{n!} f^{(n)}(0)x^n + r_n(x)
	$$
	
	其中余项
	$$
	r_n(x) = o(x^n) \quad \text{或} \quad \frac{f^{(n+1)}(\xi)}{(n+1)!} x^{n+1} \quad (\xi \text{ 介于 } 0 \text{ 与 } x \text{ 之间})
	$$
\end{corollary}
利用上面的公式，代入一些函数，可以直接得到以下函数的Maclaurin展开
\begin{proposition}
	在$x=0$的邻域内，有
	
	1.
	\[
	e^x = 1 + x + \frac{1}{2!}x^2 + \frac{1}{3!}x^3 + \cdots + \frac{1}{n!}x^n + r_n(x)
	\]
	
	推论：
	\[
	\lim_{n \to \infty} \left(1 + 1 + \frac{1}{2!} + \cdots + \frac{1}{n!}\right) = e
	\]
	
	2.
	\[
	\cos x = 1 - \frac{1}{2!}x^2 + \frac{1}{4!}x^4 - \frac{1}{6!}x^6 + \cdots + \frac{(-1)^n}{(2n)!}x^{2n} + r_{2n+1}(x)
	\]
	
	\[
	\sin x = x - \frac{1}{3!}x^3 + \frac{1}{5!}x^5 - \frac{1}{7!}x^7 + \cdots + \frac{(-1)^n}{(2n+1)!}x^{2n+1} + r_{2n+2}(x)
	\]
	
	\[
	\tan x = x + \frac{1}{3}x^3 + \frac{2}{15}x^5 + \frac{17}{315}x^7 + r_9(x)
	\]
	
	3. 对于 $f(x) = (1+x)^\alpha$, $\alpha \in \mathbb{R}$,
	
	\[\begin{aligned}
		f(x) &= (1+x)^\alpha = 1 + C_\alpha^1 x + C_\alpha^2 x^2 + \cdots + C_\alpha^{\alpha-1} x^{\alpha-1} + x^\alpha\\
		&\Rightarrow C_\alpha^0 + C_\alpha^1 x + \cdots + C_\alpha^n x^n + r_n(x) \quad (n < \alpha \text{ 时})\\
		&\Rightarrow C_\alpha^0 + C_\alpha^1 x + \cdots + C_\alpha^{\alpha-1} x^{\alpha-1} + x^\alpha \quad (n > \alpha \text{ 时})
	\end{aligned}	\]
	
	
	若$ \alpha \notin \mathbb{Z}_+$, $\forall k \in \mathbb{R}$, $C_\alpha^k \neq 0$,
	\[
	(1+x)^\alpha = 1 + \alpha x + \frac{\alpha(\alpha-1)}{2!}x^2 + \cdots + \frac{\alpha(\alpha-1)\cdots(\alpha-n+1)}{n!}x^n + r_n(x)
	\]
	
	当 $\alpha = -1$ 时，
	\[
	\frac{1}{1+x} = 1 - x + x^2 - x^3 + x^4 - \cdots + (-1)^n x^n + r_n(x)\quad(*)
	\]
	
	对于 $\alpha = \dfrac{1}{2}$，
	\[
	\sqrt{1+x} = 1 + \frac{1}{2}x - \frac{1}{2 \cdot 4}x^2 + \frac{1 \cdot 3}{2 \cdot 4 \cdot 6}x^3 - \cdots + (-1)^{n-1} \frac{(2n-3)!!}{(2n)!!}x^n + r_n(x)
	\]
	
	系数为
	\[
	\begin{cases}
		\dfrac{1}{2} & \text{当 } k=1, \\
		(-1)^{k-1} \dfrac{(2k-3)!!}{(2k)!!} & \text{当 } k > 1.
	\end{cases}
	\]
	
	对于 $\alpha = -\dfrac{1}{2}$，
	\[
	\frac{1}{\sqrt{1+x}} = 1 - \frac{1}{2}x + \frac{1 \cdot 3}{2 \cdot 4}x^2 - \frac{1 \cdot 3 \cdot 5}{2 \cdot 4 \cdot 6}x^3 + \cdots + (-1)^n \frac{(2n-1)!!}{(2n)!!}x^n + r_n(x)
	\]
	
	系数为
	\[
	(-1)^n \frac{(2n-1)!!}{(2n)!!}.
	\]
	
	对于 $(*)$，分别代入$-x,x^2,-x^2$：
	\[
	\frac{1}{1-x} = 1 + x + x^2 + x^3 + \cdots + x^n + r_n(x)
	\]
	\[
	\frac{1}{1+x^2} = 1 - x^2 + x^4 - x^6 + \cdots + (-1)^n x^{2n} + r_{2n}(x)
	\]
	\[
	\frac{1}{1-x^2} = 1 + x^2 + x^4 + \cdots + x^{2n} + r_{2n}(x)
	\]
\end{proposition}
这里介绍一条定理，用于计算一些函数的展开
\begin{proposition}
	在$x_0$的某个邻域内，$f(x)\in D^{n+2}$，则它的$n+1$次Taylor多项式的导数恰为$f'(x)$的$n$次Taylor 展开
\end{proposition}
\begin{proof}
	设
	\[
	f(x) = f(x_0) + f'(x_0)(x-x_0) + \cdots + \frac{1}{(n+1)!} f^{(n+1)}(x_0)(x-x_0)^{n+1} + r_{n+1}(x)
	\]
	
	\[
	f'(x) = P_{n+1}(x) = f'(x_0) + f''(x_0)(x-x_0) + \cdots + \frac{1}{n!} f^{(n)}(x_0)(x-x_0)^n + r'_{n+1}(x)
	\]
	
	这意味着存在 $g_n(x)$ 使得
	\[
	\lim_{x \to x_0} \frac{r'_{n+1}(x)}{g_n(x)} = 1
	\]
	
	\[
	f'(x) = f'(x_0) + f''(x_0)(x-x_0) + \cdots + \frac{1}{n!} f^{(n)}(x_0)(x-x_0)^n + g_n(x)
	\]
	
	且
	\[
	\lim_{x \to x_0} g_n(x) = 0
	\]
\end{proof}
\begin{proposition}
	1. 已知
	\[
	\frac{1}{1+x} = 1 - x + x^2 - x^3 + \cdots + (-1)^n x^n + r_n(x)
	\]
	
	设
	\[
	\ln(1+x) = a_0 + a_1 x + \cdots + a_n x^n + r_n(x)
	\]
	
	对 $\ln(1+x)$ 求导，
	\[\begin{aligned}
		(\ln(1+x))' &= a_1 + 2a_2 x + \cdots + n a_n x^{n-1} + r_{n-1}(x)\\&= \frac{1}{1+x} = 1 - x + x^2 - \cdots + (-1)^n x^n + r_{n-1}(x)
	\end{aligned}\]
	得
	\[
	\ln(1+x) = x - \frac{1}{2}x^2 + \frac{1}{3}x^3 - \cdots + (-1)^{n-1} \frac{1}{n} x^n + r_n(x)
	\]
	
	同理
	\[
	\ln(1-x) = -x - \frac{1}{2}x^2 - \frac{1}{3}x^3 - \cdots - \frac{1}{n}x^n + r_n(x)
	\]
	2. 由 $(\tan^{-1} x)' = \dfrac{1}{1+x^2} = 1 - x^2 + x^4 - x^6 + \cdots + (-1)^n x^{2n} + r_{2n}(x)$，
	得
	\[
	\tan^{-1} x = x - \frac{1}{3}x^3 + \frac{1}{5}x^5 - \frac{1}{7}x^7 + \cdots + (-1)^n \frac{1}{2n+1} x^{2n+1} + r_{2n+1}(x)
	\]
	
	类似地，
	\[
	\sin^{-1} x = x + \frac{1}{2} \cdot \frac{1}{3} x^3 + \frac{1 \cdot 3}{2 \cdot 4} \cdot \frac{1}{5} x^5 + \cdots + \frac{(2n-1)!!}{(2n)!!} \cdot \frac{1}{2n+1} x^{2n+1} + r_{2n}(x)
	\]
	
	而
	\[
	\cos^{-1} x = \frac{\pi}{2} - \sin^{-1} x
	\]
	
	3. 对于 $a^x = e^{x \ln a}$，
	\[
	a^x = 1 + x \ln a + \frac{1}{2!} (x \ln a)^2 + \cdots + \frac{1}{n!} (x \ln a)^n + r_n(x)
	\]
	
	4.
	\[
	\ln \left(\frac{1+x}{1-x}\right) = 2 \left(x + \frac{1}{3}x^3 + \cdots + \frac{1}{2n+1}x^{2n+1}\right) + r_{2n+1}(x)
	\]
	
	5.
	\[
	\cosh x = 1 + \frac{1}{2!}x^2 + \frac{1}{4!}x^4 + \cdots + \frac{1}{(2n)!}x^{2n} + r_{2n}(x)
	\]
	
	\[
	\sinh x = x + \frac{1}{3!}x^3 + \frac{1}{5!}x^5 + \cdots + \frac{1}{(2n-1)!}x^{2n-1} + r_{2n-1}(x)
	\]
\end{proposition}

